\documentclass[11pt]{article}

\usepackage[margin=25mm]{geometry}

\usepackage{amsmath,amssymb,fontspec,unicode-math}

\IfFontExistsTF{DejaVu Serif}{\setmainfont{DejaVu Serif}}{\setmainfont{Latin Modern Roman}}

\IfFontExistsTF{Latin Modern Math}{\setmathfont{Latin Modern Math}}{\setmathfont{DejaVu Math TeX Gyre}}

\usepackage{xurl,hyperref}

\setlength{\parindent}{0pt}\setlength{\parskip}{6pt}

\begin{document}

\title{Architectural and Inferential Inductive Biases for Exchangeable Sequence Modeling}

\author{ELVIA preliminary AI analysis}\date{21 October 2026}\maketitle

Preliminary assessment: Major.\par

Independent scientific review remains pending. Author adoption has not been requested.

Two Gaussian information-loss formulas are mathematically incorrect: one omits observation noise, and one mixes KL directions. These substantial derivation errors determine Major while the central conditional-information identity remains valid. The released acquisition/rollout discrepancies are Minor and do not identify the historical figure-producing implementation.

\section*{Review scope}

Latest-version, claim-focused review of the cited author-paper passages and applicable current author-release paths. The stated counterexamples and scope controls were checked analytically; no author program, old verifier, training run or benchmark was executed.

\section*{Problem 1: Bayesian linear-regression information loss omits observation noise}

Source: Appendix D.6, Example 4, p.28; common proof p.29.

Assessment: supported; impact: major

\subsection*{Source claim (paraphrase)}

Example \(4\) substitutes X Sigma X-transpose into a future-observation Gaussian determinant ratio.

\subsection*{Analysis}

Use prior variance \(2\), observation-noise variance \(2\), and repeated scalar input \(1\). After one observation, latent posterior variance is \(1\). Two future latent values therefore have singular covariance \(\begin{pmatrix}1&1\\1&1\end{pmatrix}\), so the printed noise-free determinant ratio diverges. Future observations instead have covariance \(K=\begin{pmatrix}3&1\\1&3\end{pmatrix}\), giving finite divergence \(\frac12\log(9/8)\simeq0.0588915178\). The covariance does not depend on the observed response.

\subsection*{Finding and scope}

Example4 omits observation-noise covariance and can return a divergent expression for a finite Gaussian information loss.

\subsection*{Proposed correction}

Use future-observation covariance \(K=X\Sigma X^\top+\tau^2I\) and the Gaussian ratio \(\det(\operatorname{diag}K)/\det K\), with the stated input/context expectations.

Proposed correction; independent verification and author adoption remain unestablished.

\section*{Problem 2: The GP formula mixes KL directions}

Source: Example 5, p.28; Appendix D.6.1, p.29.

Assessment: supported; impact: major

\subsection*{Source claim (paraphrase)}

Example \(5\) names forward \(\operatorname{KL}(P\|Q)\), but uses the precision of \(P\) in trace and mean terms.

\subsection*{Analysis}

For equal means, \(K_P=\begin{pmatrix}2&1\\1&2\end{pmatrix}\) and \(\Sigma_Q=2I\), the printed mixed expression gives \(0.4771743696\). Forward KL is \(\frac12\log(4/3)\simeq0.1438410362\), while reverse KL is \(0.1894922971\). Thus swapping the labels does not fix the formula. Positive observation noise makes both Gaussian laws nonsingular.

\subsection*{Finding and scope}

Example5 mixes forward and reverse Gaussian-KL terms; neither direction agrees with the printed expression.

\subsection*{Proposed correction}

For forward \(\operatorname{KL}(P\|Q)\), retain the corresponding determinant ratio and use \(\Sigma_Q^{-1}\) in trace and mean-quadratic terms. The conditional-information chain identities are not refuted.

Proposed correction; independent verification and author adoption remain unestablished.

\section*{Problem 3: The public acquisition order is fixed rather than sequential}

Source: Algorithm 3, p.25; run\_al.py:242–280,317–330.

Assessment: supported; impact: minor

\subsection*{Source claim (paraphrase)}

Algorithm 3 recomputes uncertainty after each queried label.

\subsection*{Analysis}

The script scores every candidate once before the budget loop, then repeatedly takes the first \(k\) entries of that fixed ordering. Consider two duplicate GP candidates of prior variance \(1\) and an independent candidate of variance \(0.6\), with noise \(0.01\). Querying one duplicate reduces the other's variance to \(0.01/1.01\simeq0.00990099\); sequential uncertainty selection then chooses the independent point, whereas fixed ranking need not.

\subsection*{Finding and scope}

The released active-learning path uses a fixed initial uncertainty ranking rather than the sequential rule in Algorithm3.

\subsection*{Proposed correction}

Move score computation, one query and context update into one loop, or label the existing method as fixed-ranking batch acquisition. Historical figure attribution remains unknown.

Proposed correction; independent verification and author adoption remain unestablished.

\section*{Problem 4: The rollout statistic is a final predictive mean}

Source: Algorithm 3, p.25; compute\_rollout\_uncertainty, lines 119–121.

Assessment: supported; impact: minor

\subsection*{Source claim (paraphrase)}

Algorithm 3 averages all J generated samples within each path before comparing paths.

\subsection*{Analysis}

The release records the final predictive mean, rather than averaging all \(J\) generated samples along the path. For latent variance \(v\) and noise variance \(\tau^2\), the path-mean variance is \(v+\tau^2/J\), while the final-predictive-mean variance is \(v-v\tau^2/[\tau^2+(J-1)v]\). At \(v=\tau^2=1,J=1\), they are \(2\) and zero; at \(J=20\), they are \(1.05\) and \(0.95\). Both approach \(v\), so the alternative statistic is not automatically useless.

\subsection*{Finding and scope}

The release uses a final predictive mean instead of the declared within-path sample mean. Its asymptotic agreement does not make the finite statistics equal.

\subsection*{Proposed correction}

Compute the stated path-sample mean or explicitly define and analyze the alternate finite-path estimator. Corrected learning curves are not asserted.

Proposed correction; independent verification and author adoption remain unestablished.

\section*{Limits}

Independent scientific adjudication remains pending.\par

No author code, model training or full benchmark evaluation was executed in this revision.\par

Version-specific source defects do not establish which code produced the published experiments.\par

Corrected benchmark scores, empirical effect sizes and author adoption remain unverified unless explicitly stated.\par

No new author-code execution, model spending, training, author contact or complete benchmark replication.\par

Historical numerical results are retained evidence rather than current reruns.\par

Local counterexamples do not invalidate every empirical result; missing historical run attribution remains unresolved.\par

\section*{Source identity}

NeurIPS 2025 official proceedings PDF (current publisher bytes, checked 2026-10-21)

\url{https://proceedings.neurips.cc/paper_files/paper/2025/file/e31fc00e6acb1d09184f1489e766900d-Paper-Conference.pdf}

Reviewed PDF SHA-256: \nolinkurl{cb6875508f7d576ee9735e6bbe947aec90862c02e6eeef8c0270b96aa61b1486}

This is the same research work as the predecessor, matched by title, author identity and official publisher channel. Its freshly retrieved exact PDF bytes and real retrieval time are recorded. Current target arXiv metadata and the relevant available author-version passages were checked separately; background works were excluded. Original report/source history is unchanged.

\end{document}
